\chapter{Cohorts}
\label{chap:cohorts}

A \emph{cohort} is a named set of patients, found across every database by criteria, and kept.
The \ui{Cohorts} tab, the first of the header and the one the page opens on, searches as the
criteria change, shows why each patient matched, and saves all the matches (or the ticked ones)
as a new cohort, or adds them to one. A cohort keeps its members until someone refreshes it; its
definition records its criteria, where its members come from, and every patient added or
removed manually.

\screenshot{cohorts.png}{The Cohorts tab: an expression, the patients it finds and why each
matched.}{fig:cohorts-builder}

\section{The Cohorts tab}
\label{sec:cohorts-tab}

On the left, the tab lists the cohorts that the user may see: the user's own, those shared with
the user, and the demo's. Each shows its number of patients and whose it is (``shared'' when the
user shared it, ``from \dots'' with its owner's name when it was shared with the user, ``the
demo''). A subcohort is listed under its parent. A click opens the cohort's page
(\cref{sec:cohorts-page}).

Above the list, \ui{New / Edit cohort} starts a new cohort. On a cohort's page it is a menu with
\ui{New cohort}, \ui{Edit \dots} (when the user may edit the cohort), \ui{Clone \dots}, \ui{New
subcohort of \dots}, \ui{Merge \dots{} with others\dots}, and \ui{Delete \dots} (for its owner,
and for the demo's cohort, any administrator). \key{Esc} closes the menu.

\section{Building the criteria}
\label{sec:cohorts-builder}

The criteria builder says who is \emph{included} and who is \emph{excluded}:
\begin{itemize}
\item \textbf{Include patients who meet any group (and all of a group's conditions).} The
  conditions of a group are joined by \emph{and}, and the groups by \emph{or}. Without any
  inclusion condition, every patient is included.
\item \textbf{Exclude those who ever had.} An exclusion takes out every included patient who
  \emph{ever} met it, in their whole history. The results list who was excluded, and why.
\end{itemize}

A new cohort starts with one expression line, where criteria are written as text
(\cref{sec:cohorts-expressions}). The following buttons add lines:
\begin{description}
\item[\ui{+ and\dots}] A condition in the same group (a \ui{Lab or vital sign} line, whose kind
  can then be changed).
\item[\ui{+ and an expression}] An expression in the same group.
\item[\ui{+ or another group}, \ui{+ Add a condition}] A new group (the latter when there is
  none yet).
\item[\ui{+ Write an expression}] A new group holding an expression.
\item[\ui{+ Add an exclusion}] An exclusion (a \ui{Diagnosis} line, whose kind can be changed).
\end{description}
The $\times$ at the end of a line removes it. Each inclusion line begins with \ui{with} or
\ui{without} (\ui{without} selects the patients who do not meet the condition), then the kind of
condition (\cref{tab:cohorts-kinds}), and then its fields.

\begin{xltabular}{\linewidth}{@{}P{3.4cm}P{4.2cm}L@{}}
\caption{The kinds of condition of the builder.}
\label{tab:cohorts-kinds}\\
\toprule
{\tabhead Kind} & {\tabhead Fields} & {\tabhead A patient meets it when} \\
\midrule
\endfirsthead
\toprule
{\tabhead Kind} & {\tabhead Fields} & {\tabhead A patient meets it when} \\
\midrule
\endhead
\bottomrule
\endfoot
\ui{Expression} & The text & They meet the expression (\cref{sec:cohorts-expressions}) \\
\ui{Lab or vital sign} & Name, value, \ui{latest only}, the database (\ui{any database}),
  dates & A result or vital sign of that name has the value, within the dates; with a database
  chosen, of exactly that name in it, else of any name containing it \\
\ui{A database's column} & The database, the column, \ui{value} or \ui{contains}, the value or
  the text, \ui{latest only} & A row of that database about them has the value (or the text) in
  that column \\
\ui{Diagnosis} & Code or name & A condition's code starts with it (\code{I21}), or its name
  contains it (\code{sepsis}) \\
\ui{Medication} & Name, dates & A drug or transfusion whose name contains it was given within
  the dates \\
\ui{Has data} & A category & They have data of that category (\code{ICU}, \code{Holter},
  \code{EEG}, \code{Sleep}, \code{ECG} and so on) \\
\ui{Encounter} & Unit, ICU (\ui{any}, \ui{with an ICU stay}, \ui{without an ICU stay}), dates
  & An encounter overlaps the dates, in a unit whose name contains the text, with or without an
  ICU stay \\
\ui{Sex} & \code{F} or \code{M} & Their sex is that \\
\ui{Age} & Value & Their age now (in years) has the value; the option is disabled (its title:
  ``Age comes from the date of birth: it needs the PHI permission, out of Privacy Mode'') unless
  the user has the PHI permission and is out of Privacy Mode \\
\ui{Text in notes and reports} & The text (at least 3 letters) & An item's text contains it \\
\ui{In a database} & The database & The database holds data about them \\
\end{xltabular}

\paragraph{Values.} A value is a comparison and a number. The comparisons are \ui{between}
(\ui{from} and \ui{to}, inclusive, either of which may be left empty), \ui{>}, \ui{$\geq$},
\ui{<}, \ui{$\leq$}, \ui{=} and \ui{$\neq$}. A unit may follow the number (the field offers
\code{ms}, \code{s}, \code{\textmu V}, \code{mV}, \code{bpm}, \code{\textdegree}, \code{mmol/L}
and \code{mg/dL}), and it is converted to the data's unit (\cref{sec:cohorts-units}). With
\ui{latest only}, only each patient's latest value of that exact name counts (\code{latest QT}
never reads a QTc), across the catalog and the connected databases; otherwise any value, at any
time, counts.

\paragraph{Dates.} \ui{from} and \ui{to} are days in the catalog's time zone, \ui{to} included.
In Privacy Mode, the dates shown are moved back, and the dates of a condition (an encounter's,
a result's, a drug's) are dates as shown: each patient's own dates are compared with them, so the
cohort finds the patients whose shown dates fall in the range, and the criteria keep that meaning
for every user who refreshes them (\cref{sec:protecting-dates}). An encounter without a discharge
time counts as continuing until now.

\begin{phibox}
\ui{Age} comes from the date of birth. The server answers it only for users with the PHI
permission who switched Privacy Mode off (``age comes from the date of birth: it needs the PHI
permission''), and the builder disables the \ui{Age} option for everyone else. \ui{Text in notes
and reports} searches the text of notes only for users who may read it (\cref{chap:patients});
otherwise notes are left out. A column that identifies (MRN, names, date of birth, encounter
number), or that is not mapped and is not a measurement, can be used only with the PHI
permission, out of Privacy Mode. Searches made with the PHI permission that use age, text or
expressions are written to the audit log (\code{phi\_search}), and so are criteria searched in
Privacy Mode (\code{privacy\_search}).
\end{phibox}

\subsection{The results}
\label{sec:cohorts-results}

The search runs as the criteria change. An expression is checked as it is typed, and its
meaning, in words, appears beside it (or the error; \cref{sec:cohorts-errors}). Above the
results, the tab says ``\emph{N} patients match'', with the number excluded and the number
ticked. The table lists the patients found, a hundred to a page, with \ui{Patient},
\ui{Patient ID}, \ui{Sex}, \ui{Age}, \ui{Data} (the colors of their categories), and \ui{Why}:
up to three reasons for each patient (``QTc 482 ms (ecg\_index) on 2024-01-03'', ``has ICU'',
``no Holter''). A click on a patient's name opens the patient's page. Below the table, a fold
lists the patients excluded, with \ui{Excluded for}.

\paragraph{Saving.} The boxes at the start of the rows tick patients, and the box in the header
ticks the whole page. Below the criteria, two choices save the result:
\begin{itemize}
\item \ui{Save all \emph{N} as a new cohort} (or \ui{Save the \emph{K} ticked as a new cohort}),
  after a name is typed in the field \ui{New cohort's name};
\item \ui{or add them to} a cohort that the user may change, followed by \ui{Add}. The patients
  whom the cohort's definition does not select are recorded in it as added manually.
\end{itemize}
The new cohort's definition is the criteria as written, and its members are the patients saved.

\section{The expression language}
\label{sec:cohorts-expressions}

Any line may be an expression: criteria written as text, parsed by the server
(\file{medorama/src/expr.rs}) into the same structured query as the builder's lines. For
example:
\begin{lstlisting}[style=text]
(age > 18 AND age < 85 AND female) OR (age > 18 AND age < 65 AND male AND QT > 400ms)
\end{lstlisting}

\subsection{Grammar}
\label{sec:cohorts-grammar}

\begin{lstlisting}[style=text]
expression  :=  or
or          :=  and  { OR and }
and         :=  factor  { [AND] factor }       two conditions side by side mean both
factor      :=  NOT factor  |  "(" or ")"  |  condition
condition   :=  sex | age | has | diagnosis | medication | text | in | latest | name-test
\end{lstlisting}

\begin{itemize}
\item \code{AND} binds before \code{OR}: \code{a OR b AND c} is \code{a OR (b AND c)}.
\item The words \code{AND}, \code{OR}, \code{NOT} are read in any case (\code{and}, \code{Or});
  \code{\&\&}, \code{||} and \code{!} mean the same.
\item Two conditions written side by side, with nothing between them, mean both:
  \code{female has ecg} is \code{female AND has ecg}.
\item \code{NOT} applies to the factor that follows it: a condition or a bracketed group. It
  keeps the patients who do \emph{not} meet it, including those who have no such data at all:
  \code{NOT QT > 400ms} keeps the patients without any QT above 400\,ms, and those without a QT.
\end{itemize}

\paragraph{Tokens.}
\begin{itemize}
\item \textbf{Words} are made of letters, digits and \code{\_}, and may hold \code{.}, \code{-},
  \code{/}, \code{\%}, \code{\textmu} and \code{\textdegree} after their first character
  (\code{ecg\_index.QT}, \code{mmol/L}).
\item \textbf{Quoted text}, between \code{"} and \code{"}, \code{'} and \code{'}, or typographic
  double quotes, is one token, spaces included: \code{"R axis"}, \code{"pulmonary edema"}.
  Quoted text is never a keyword.
\item \textbf{Numbers}: digits, with a sign, a decimal point and an exponent
  (\code{-30}, \code{0.4}, \code{1e3}). A unit may be written against the number:
  \code{400ms}, \code{1.2mV}.
\item \textbf{Comparisons}: \code{>}, \code{<}, \code{>=}, \code{<=}, \code{=} (or \code{==}),
  \code{!=}, and the characters $\geq$, $\leq$, $\neq$. \textbf{Text}: \code{contains} (or
  \code{mentions\_text}, or \code{\textasciitilde}).
\end{itemize}

\subsection{Conditions}
\label{sec:cohorts-conditions}

\Cref{tab:cohorts-conditions} lists every condition the language knows. Keywords are read in any
case. Where a condition takes a \emph{text}, it is a quoted text or a single word (or a number),
optionally after \code{=} or \code{contains}: \code{dx I21}, \code{dx = I21},
\code{diagnosis "septic shock"}; several words need quotes.

\begin{xltabular}{\linewidth}{@{}P{5.6cm}L@{}}
\caption{The conditions of the expression language.}
\label{tab:cohorts-conditions}\\
\toprule
{\tabhead Written} & {\tabhead A patient meets it when} \\
\midrule
\endfirsthead
\toprule
{\tabhead Written} & {\tabhead A patient meets it when} \\
\midrule
\endhead
\bottomrule
\endfoot
\code{female}, \code{females}, \code{women}, \code{woman}, \code{F} & Their sex is F \\
\code{male}, \code{males}, \code{men}, \code{man}, \code{M} & Their sex is M \\
\code{sex} \emph{text}, \code{gender} \emph{text} & Their sex is the first letter of the text
  (\code{sex = female}, \code{gender M}) \\
\code{age} \emph{comparison number} [\emph{unit}] & Their age now, in years, compares so
  (\code{age >= 65}, \code{age > 18 years}); the unit, if any, is \code{y}, \code{yr},
  \code{yrs}, \code{year} or \code{years}, in any case. Needs the PHI permission, out of Privacy
  Mode \\
\code{has} \emph{category}, or a category's word alone & They have data of that category:
  \code{has ecg}, \code{icu}, \code{holter}, \code{sleep}, \code{eeg}, \code{echo} and so on
  (any word of a category, \cref{tab:patients-categories}) \\
\code{diagnosis}, \code{dx}, \code{condition} \emph{text} & A condition's code starts with the
  text, or its name contains it: \code{dx I21}, \code{diagnosis sepsis} \\
\code{medication}, \code{medications}, \code{med}, \code{drug} \emph{text} & A drug or
  transfusion whose name contains the text: \code{medication heparin} \\
\code{text}, \code{mentions}, \code{notes} \emph{text} & An item's text contains it (at least 3
  characters): \code{text "pulmonary edema"}; the text of notes only for users who may read
  it \\
\code{in} \emph{database} & The database holds data about them: \code{in ecg\_index},
  \code{in "ecg\_index"} \\
\code{latest} (or \code{last}) \emph{name-test} & Only each patient's latest value of a
  column or a measurement of that exact name counts (\code{latest QTc >= 470}; \code{latest QT}
  never reads a QTc), across the catalog and the connected databases; without it, any value,
  ever \\
\emph{name comparison number} [\emph{unit}] & A database's column, or results of that name,
  compares so, in its unit: \code{QT > 400ms}, \code{"R axis" < -30},
  \code{ecg\_index.QT > 0.4 s}, \code{labs.creatinine > 1.5} \\
\emph{name} \code{contains} \emph{text}, \emph{name} \code{\textasciitilde} \emph{text} & A
  column's text contains it, ignoring case:
  \code{ecg\_index.interpretation contains "fibrillation"} \\
\emph{name} \code{=} \emph{text}, \emph{name} \code{!=} \emph{text} & A column's text is (or is
  not) exactly that, ignoring case: \code{ecg\_index.interpretation = "Sinus rhythm"} \\
\emph{database}\code{.}\emph{name} \dots & The same, in the named database (connected
  databases included: \code{demo\_omop.lactate > 4}) \\
\end{xltabular}

A word that is neither a keyword nor a category must be followed by a comparison or
\code{contains}: \code{zzz} alone is the error ``what is "zzz"? (a name needs a comparison:
zzz > 1)''.

\subsection{Names: columns and results}
\label{sec:cohorts-names}

A \emph{name} is written alone (\code{QT}), in quotes when it has spaces (\code{"R axis"}), or
after a database's name and a dot (\code{ecg\_index.QT}, \code{ecg\_index."R axis"}). A
database is named as the Databases tab shows it (its file's name without \file{.csv},
\file{.tsv}, \file{.txt}, \file{.xlsx} or \file{.xls}), ignoring case, or by its number
(\code{\#12}); a beginning that only one database's name starts with is enough.

\paragraph{Columns are never merged across databases.} A name refers to, in this order:
\begin{enumerate}
\item \textbf{a column} of a database: one whose header is the name, or whose header without its
  unit is (\code{QT} for \code{QT (ms)}), or the one column whose name begins with those words
  (\code{LAD} for \code{LAD stenosis (\%)}). When several databases have such a column (an ECG
  database and a Holter database both with \code{QT}), the expression must say which, and the
  error lists them: ``"QT" is a column of several databases (\dots): write which, e.g.
  ecg\_index.QT'';
\item else \textbf{results of that name} (items with a numeric value, such as a laboratory
  result or a vital sign), ignoring case, in one database;
\item else \textbf{a connected database's measurement}: one of that name, else the one name that
  contains it. OMOP's concept names are long: when several contain the word, the error lists
  them (``"lactate" names 3 measurements in omop: \dots; write which'') and the one meant is
  written in quotes: \code{omop."Lactate [Moles/volume] in Blood" > 4}.
\end{enumerate}
When the results of that name are in several databases, the database must be named too.
Nothing found is the error ``nothing is named \dots''.

\subsection{Units}
\label{sec:cohorts-units}

A comparison's number is compared in the data's unit, converted from the unit written after it:
\code{QT > 400ms} and \code{QT > 0.4 s} are the same condition. The units written after a number
are \code{ms}, \code{s}, \code{sec}, \code{min}, \code{h}, \code{uV}, \code{\textmu V},
\code{mV}, \code{V}, \code{bpm}, \code{\%}, \code{kg}, \code{g}, \code{mg}, \code{cm},
\code{mm}, \code{m}, \code{\textdegree}, \code{deg}, \code{mmHg}, \code{mmol/L}, \code{mg/dL},
\code{g/dL}, \code{ng/L}, \code{ng/mL}, and any unit written with a \code{/}, in any case.

\begin{table}[htbp]
\caption{The units converted into one another.}
\label{tab:cohorts-units}
\begin{tabularx}{\linewidth}{@{}P{2.4cm}L@{}}
\toprule
{\tabhead Quantity} & {\tabhead Units} \\
\midrule
Time & \code{s} (\code{sec}, \code{secs}, \code{second}, \code{seconds}), \code{ms}
  (\code{msec}), \code{us}, \code{\textmu s} (\code{usec}), \code{min} (\code{minutes}),
  \code{h} (\code{hr}, \code{hours}) \\
Voltage & \code{V}, \code{mV}, \code{\textmu V} (\code{uV}) \\
Rate & \code{bpm}, \code{/min}, \code{beats/min} \\
Angle & \code{\textdegree}, \code{deg}, \code{degrees} \\
Mass & \code{kg}, \code{g}, \code{mg} \\
Length & \code{m}, \code{cm}, \code{mm} \\
\bottomrule
\end{tabularx}
\end{table}

\begin{itemize}
\item Units of the same quantity are converted: \code{0.15 mV} is \code{150 \textmu V}.
\item The same unit, written in another case, compares as it is: \code{mmol/L} and
  \code{mmol/l}.
\item A number written without a unit is compared as it is, and only when the data hold one
  unit: when they hold several, write the unit (\code{labs.creatinine > 1.5}). An age is in
  years, with or without a unit of years (\code{age > 18}, \code{age > 18 years}). After a
  number, \code{m} is read as \emph{male}, not as metres.
\item Values in a unit that does not compare with the one written (mg/dL against
  \code{glucose > 7 mmol/L}), or values without a unit, are left out, and the search says how
  many. When none of the values compares, the condition is an error: units of different
  quantities (\code{ms} and \code{mV}), or two different units not in \cref{tab:cohorts-units}
  (\code{mmol/L} and \code{mg/dL} are not converted): ``QT is in ms: 5 mV does not compare with
  it''.
\end{itemize}
\textbf{Saved criteria keep their meaning.} When a cohort is saved, each expression is kept
together with what its names meant at that time, such as which database's \code{QT} column. A
database imported later that would make a name ambiguous, or give it another meaning, does not
change the cohort: a search of its criteria says so, and the cohort keeps the meaning it was
saved with.
A column's unit is the one its header gives in brackets (\code{QT (ms)}, \code{Lactate
[mmol/L]}), or, for a few known ECG measurements written without one, the usual unit:
\code{Vent. rate}, \code{Ventricular rate}, \code{Atrial rate}, \code{Heart rate} and \code{HR}
in bpm; \code{PR}, \code{QRS}, \code{QT}, \code{QTc}, \code{QTcB}, \code{QTcF} and \code{RR} in
ms; \code{P axis}, \code{R axis}, \code{QRS axis} and \code{T axis} in degrees. A result's unit
is the one its row gives.

\subsection{Errors}
\label{sec:cohorts-errors}

An expression that cannot be read shows its error beside it, with the position of the problem
(``at 12: \dots''), and the search does not run until it is fixed. The errors include:
\begin{itemize}
\item ``the expression is empty''; ``a quote is not closed''; ``a ')' is missing'';
\item ``a condition is expected here''; ``a condition is missing'';
\item ``a comparison (>, <, >=, <=, =, !=) is expected''; ``a number is expected'';
\item ``what is "\dots"? (a name needs a comparison: \dots{} > 1)'';
\item ``this is not understood (an AND or OR missing?)'';
\item ``no category "\dots" (ecg, icu, holter, imaging \dots)'';
\item ``age is in years''; ``latest applies to a measurement or a column'';
\item ``text is compared with = or contains'';
\item ``no database "\dots" (databases: \dots)''; ``"\dots" names several databases: \dots'';
\item the errors of names (\cref{sec:cohorts-names}) and units (\cref{sec:cohorts-units}).
\end{itemize}

\subsection{Examples}
\label{sec:cohorts-examples}

On the demo (\cref{chap:quickstart}):
\begin{lstlisting}[style=text]
labs.creatinine > 1.5 OR has icu              a result of the database labs, or an ICU stay
QTc > 470ms AND female                         a column of ecg_index, in its unit (ms)
latest ecg_index.QTc >= 470                    only each patient's latest ECG counts
ecg_index.interpretation contains "fibrillation"
"R axis" < -30                                 a column name with a space, in quotes
has holter AND NOT has icu
has sleep OR has eeg                           a sleep study or an EEG recording
dx I21 OR diagnosis "myocardial infarction"
medication heparin AND in icu_index
text "pulmonary edema"
(female OR age > 80) AND has echo              age needs the PHI permission
\end{lstlisting}

\section{A cohort's page}
\label{sec:cohorts-page}

At the top, a cohort's page shows its name, its number of patients, whose it is (``yours'';
``\dots's, shared with you to read'' or ``to edit''; ``the demo, for every user''), when it was
made, when its members were found, and the version of each connected database at that time. Its
description follows, and then:
\begin{itemize}
\item the buttons \ui{Edit} (when the user may edit the cohort), \ui{Clone\dots},
  \ui{Subcohort\dots}, \ui{Merge\dots}, \ui{Share\dots} (for its owner) and \ui{Delete\dots}
  (for its owner, and for the demo's cohort, an administrator);
\item its \textbf{definition} (\cref{sec:cohorts-definition});
\item when the catalog has new data since the members were found, the notice ``New data since
  the members were found'', with \ui{See who would join\dots};
\item \ui{Refresh the members\dots} (\cref{sec:cohorts-refresh}), \ui{Open in Patient Catalog}
  (its members as the patient list, \cref{chap:patients}), \ui{Run actions\dots} (on its
  members' recordings, \cref{chap:actions}), and the exports
  (\cref{sec:cohorts-export});
\item the \textbf{members} (\cref{sec:cohorts-members});
\item its \textbf{History} (\cref{sec:cohorts-history}).
\end{itemize}

\subsection{The definition}
\label{sec:cohorts-definition}

A cohort's definition is shown in words:
\begin{description}
\item[\ui{Made}] How the cohort was made (cloned from another cohort, the demo, and so on).
\item[\ui{From}] Where its members come from, when not from the whole catalog: ``Subcohort of
  \dots'', ``In \dots{} or \dots'' (a union), ``In \dots{} and \dots'' (an intersection), or ``In
  \dots{} but not in \dots'' (a difference). Each cohort named is a link, and a deleted one is
  marked ``(deleted: its last members are kept)''.
\item[\ui{Criteria}] What includes a patient, and each exclusion (``never: \dots''). Without
  criteria, it reads ``none: every member of its source'' or ``none: its members are added
  manually''.
\item[\ui{Added manually}, \ui{Removed manually}] The patients added or removed by hand,
  identified by patient ID, never by name or MRN. A long list folds after the first 40 (\ui{and
  \dots{} more}).
\end{description}

\subsection{Members}
\label{sec:cohorts-members}

The members' table shows each patient as the user's permissions allow (identifiers only with
the PHI permission, pseudonyms in Privacy Mode), and how and when the patient was added:
\ui{by the criteria}, \ui{by the definition} (for a cohort with a source), or \ui{added
manually}. A user who may edit the cohort can do the following:
\begin{itemize}
\item Remove a member with the $\times$ at the end of its row (``Remove from the cohort
  (recorded in its definition)''). The patient is then removed manually, and stays out when the
  cohort is refreshed.
\item Add patients by search. The box above the members runs the Patient Catalog's search
  (\cref{chap:patients}) and lists the patients found, with members marked as such. \ui{Add the
  \emph{K} ticked} or \ui{Add all \emph{N} found} adds them, and \ui{Clear} empties the search.
  Patients whom the definition does not select are recorded as added manually, and adding a
  patient who was removed manually ends that removal.
\end{itemize}

\section{Keeping and refreshing}
\label{sec:cohorts-refresh}

A cohort keeps its members as a list, a snapshot: new data change nothing until someone
refreshes it. Its page says when the catalog has new data since its members were found.
\ui{Refresh the members\dots} (or \ui{See who would join\dots}) first shows what the definition
selects now (\ui{What the definition selects now}):
\begin{itemize}
\item the new patients who would join, each ticked (untick those not wanted);
\item in a fold, the members whom the definition no longer selects, who would leave.
\end{itemize}
Members added manually stay, and members removed manually stay out. Three buttons then decide:
\begin{description}
\item[\ui{Add the \emph{N} new patients}] (or \ui{Add \emph{K} chosen patients}) Adds the
  ticked patients, and nobody leaves.
\item[\ui{Apply every change (\emph{N} join, \emph{M} leave)}] Applies the changes shown: the new
  patients ticked join, and the members shown as leaving leave. Changes that appeared since the
  preview are not applied, and the message says how many: refresh again to see them.
\item[\ui{Cancel}] Changes nothing.
\end{description}
When nothing would change, the page says ``The members are up to date''. A refresh reads the data
as the criteria were written: criteria on notes' text, ages, columns or expressions select other
patients for a user who reads the data differently (in Privacy Mode instead of with the PHI
permission, for example), so such a refresh is refused, and the message says how the criteria
were written. A subcohort or a merge follows the cohorts it is made from only while the user may
still read them; when one is no longer shared with the user, a refresh is refused. A refresh records the
catalog's data version and each connected database's version at that time; its entry in the
cohort's history names the connected databases that had new data.

\section{Editing}
\label{sec:cohorts-edit}

\ui{Edit} opens the builder with the cohort's \ui{Name}, \ui{Description} and criteria (and,
for a subcohort or a merge, where its members come from, set when it was made). The results show
the patients the edited criteria select. \ui{Save the changes} saves the name, the description
and the criteria; the criteria are then applied by a refresh, shown first: the cohort's page
opens with who would join and who would leave. \ui{Cancel} changes nothing. Members are changed
on the cohort's page, not in the builder.

Criteria made elsewhere (through the API, \cref{chap:http-api}) that the builder cannot draw
are shown read-only and kept as they are:
\begin{itemize}
\item inclusion criteria that are not groups of conditions are shown in words (``Criteria made
  elsewhere, kept as they are''), until \ui{Replace them};
\item an exclusion that is a group is kept as a read-only line (``A group made elsewhere, kept as
  it is (not editable here)''), and saved unchanged; the $\times$ at its end removes it.
\end{itemize}

\section{Subcohorts, merges and clones}
\label{sec:cohorts-derived}

\paragraph{A subcohort} (\ui{Subcohort\dots}, or \ui{New subcohort of \dots}) is made with the
builder, which then searches among the parent's members only. \ui{Save \dots{} as a subcohort
of \dots} saves it. It follows its parent when it is refreshed: its members are chosen among the
parent's members of that time.

\paragraph{A merge} (\ui{Merge\dots}) makes a new cohort from this one and others, ticked in a
list:
\begin{description}
\item[\ui{Union}] The patients in this cohort or in any of the others.
\item[\ui{Intersection}] The patients in this cohort and in all of the others.
\item[\ui{Difference}] The patients in this cohort and in none of the others.
\end{description}
\ui{Make the cohort}, after a name is typed, makes it. The merged cohorts are left as they are,
and the new cohort follows them when it is refreshed.

\paragraph{A clone} (\ui{Clone\dots}) is a cohort of the user's own, with the same definition
and members (and following the same cohorts when refreshed); its name is first proposed as
``\dots{} (copy)''. The cohort cloned does not change. Any cohort the user may see can be
cloned.

When a cohort that others are made from is deleted, they keep its last members and follow
those from then on.

\section{Sharing}
\label{sec:cohorts-share}

A cohort is private to its owner unless shared. \ui{Share\dots} (its owner only) lists the
other users, each with \ui{May see it} and \ui{May edit it}; \ui{Save} applies the choice.
Those who may edit it change its name, description, criteria and members, and refresh it; only
its owner shares and deletes it. Each user sees the members as their own permissions allow:
identifiers only with the PHI permission, out of Privacy Mode.

\section{Deleting}
\label{sec:cohorts-delete}

\ui{Delete\dots} (its owner only) says first what happens: the patients in the cohort stay in
the catalog, and the cohorts made from it keep its current members. \ui{Delete} then deletes it;
it cannot be undone.

\begin{note}[The cohort DEMO]
The cohort \code{DEMO} exists when Medorama was installed with its demo. It is the one cohort
whose deletion also removes its patients. Every user sees it, read-only, and any administrator
may delete it. \ui{Delete\dots} first says what goes (its patients, encounters, bed stays, items
and recordings, and its OMOP database \code{demo\_omop}), and \ui{Delete the demo} then removes
all of it from the catalog (\cref{chap:quickstart}).
\end{note}

\section{History}
\label{sec:cohorts-history}

Each change to a cohort (made, renamed, criteria changed, members added or removed, refreshed,
shared, cloned) is in its \ui{History}, at the bottom of its page (\ui{When}, \ui{Who},
\ui{What}; ``command line'' for changes made there), and in the audit log
(\cref{chap:administration}).

\section{Exports}
\label{sec:cohorts-export}

The cohort's page exports the cohort as files to download (\cref{tab:cohorts-exports}). In
Privacy Mode, the label reads \ui{Export (Privacy Mode)}, and the exports pass the privacy gate.
They have the same columns, with identifiers as placeholders, patients named by their
pseudonyms, dates moved back and ages above 89 as \code{90+}, and their file names end in
\code{-privacy}.

\begin{table}[htbp]
\caption{A cohort's exports.}
\label{tab:cohorts-exports}
\begin{tabularx}{\linewidth}{@{}P{3.0cm}P{4.4cm}L@{}}
\toprule
{\tabhead Button} & {\tabhead File} & {\tabhead Holds} \\
\midrule
\ui{Members (CSV)} & \file{cohort-NAME.csv} & The members, one row per patient; for users with
  the PHI permission, their identifiers too (MRN, names, date of birth, age) \\
\ui{Definition (JSON)} & \file{cohort-NAME-criteria.json} & The criteria and the source as
  stored, and the whole definition in words, with the patients added and removed manually \\
\ui{Members' data (CSV)} & \file{cohort-NAME-data.csv} & Every database's rows about the
  members, one cell per line: patient, database, line, column, value \\
\ui{Excel} & \file{cohort-NAME.xlsx} & The members, the definition, and a sheet per database \\
\bottomrule
\end{tabularx}
\end{table}

The columns exported are those the user may see (\cref{chap:databases-tab}). Every member is
exported, however many. Times are written in UTC, marked \code{Z} (ISO 8601). A cell that a
spreadsheet would run as a formula (text starting with \code{=}, \code{+}, \code{-} or
\code{@}) is written after a \code{'}, so that it opens as text.

\section{Connected databases in cohorts}
\label{sec:cohorts-connected}

Criteria are answered by every source that holds such data: the catalog, and each connected
database (an OMOP instance, a warehouse; \cref{chap:connecting}). A condition, a drug, a text, a
stay or a result found in any of them counts; sex, age and ``data in that database'' are
answered by the catalog for every linked person. A connected database is named in criteria by
the name it was connected with (\code{omop.lactate > 4}). Its persons are linked to the
catalog's patients as their data are found, so that a person in several sources is one patient.
A cohort records the version of each connected database when its members were found, shown on
its page, and a refresh says which have new data.

\section{Criteria as data}
\label{sec:cohorts-json}

Behind the builder, the criteria are a structured query, which the API takes and returns
(\cref{chap:http-api}), and which another program can write:
\begin{lstlisting}[language=JSON]
{"version": 1,
 "include": {"op": "or", "of": [
     {"kind": "expression", "text": "QTc > 470ms AND female"},
     {"op": "and", "of": [{"kind": "has", "category": "icu"},
                          {"kind": "measurement", "name": "Lactate", "cmp": ">", "value": 4}]}]},
 "exclude": [{"kind": "diagnosis", "text": "I21"}]}
\end{lstlisting}
An expression is a criterion or a group (\code{"op"}: \code{"and"}, \code{"or"} or
\code{"not"}, \code{"of"}: a list); a criterion's \code{kind} is one of those of
\cref{tab:cohorts-kinds} (\code{expression}, \code{measurement}, \code{column},
\code{diagnosis}, \code{medication}, \code{has}, \code{encounter}, \code{sex}, \code{age},
\code{text}, \code{database}), with \code{not} to invert it. A patient is in when they meet
\code{include} (none: everyone) and no \code{exclude}, looked for in their whole history. The
format: \cref{chap:formats-reference}.
